\chapter{Type theoretical formulation of self-adaptive robotic systems}
\label{type}

\begin{abstract}
Sample Abstract.
\end{abstract}

\blfootnote{This chapter is a verbatim copy of the paper~\cite{}.
\begin{itemize}
    \item Type theoretical formulation of self-adaptive robotic systems. In preparation for a short paper for SEAMS 2025.
\end{itemize}
}

\section{Abstract}

\section{Introduction}

\textbf{Research question:} 
\begin{itemize}
    \item How to formally define the self-adaptation problem in robotic systems? 
\end{itemize}


common ``parent'' definition of self-adaptive robotic systems that can be specialized for different techniques

clear definition of the SARS problem to prevent ambiguities and facilitate the comparison of different techniques

combination of robotics architecture with self-adaptive robotic systems

composable theory, so one can take only parts

extensible theory, one can increment it

showcases for metacontrol, ROSA, software product lines?

% check garlan paper on 2nd level of adaptation as a reference for "standard" description of adaptation

\section{MAPE-K formulation}

Monitor:
\begin{gather}
    SensorInput : Type \\
    MonitorOutput: Type \\
    s: SensorInput \vdash monitor(s) : MonitorOutput \\
\end{gather}

Analyze:
\begin{gather}
    AnalyzeOutput: Type \\
    mo : MonitorOutput \vdash analyze(mo) : \{ ao>AnalyzeOutput \}\\
\end{gather}


Plan:
\begin{gather}
    PlanOutput: Type \\
    ao: AnalyzeOutput \vdash plan(ao) : PlanOutput 
\end{gather}

Execute:
\begin{gather}
    ExecuteOutput : Type \\
    ao : AnalyzeOutput \vdash execute(rp) : \{ eo>ExecuteOutput \}, 
\end{gather}

One possible definition of the MAPE-K loop:
\begin{lstlisting}
    while true:
        s = sensor_inputs()
        mo = monitor(s)
        ao = analyze(mo)
        rp = plan(ao)
        execute(rp)
\end{lstlisting}

\section{Robotics architecture formulation}

\textbf{Deliberative layer:} Devise task plan. \textbf{Input:} robot's goal. \textbf{Output:} task sequence.

\textbf{Executive layer:} Task execution. \textbf{Input: task}. \textbf{Output:} action sequence.

\textbf{Functional layer (skill layer/behavioral layer):} Realize the robot's functions as components

task planner:

components ...

\begin{gather}
    Action: Type \\
    Mode: Type \\
    Component: Type \\
    Condition: Type 
\end{gather}

% \section{Definitions}

% \section{Robot architecture formulation}

\section{Self-adaptative robotic systems formulation}

The skills the robot has are defined as $S = \{S_1, ..., S_n\}$. The actions the robot can perform are defined as $A = {A_1, ..., A_m}$, where each action $A_i$ is defined as a set of skills required to perform it $A_i=\{S_j,..., S_p\}$. The possible modes (i.e., active components or component's parameterizations) of a skill layer with $n$ skills are defined as $\theta = \{\theta_1, ..., \theta_n\}$, where $\theta_i$ is the set of possible modes for a skill $S_i$. The set of conditions for the selection of modes in the skill layer is defined as $\sigma = \{\sigma_1, ..., \sigma_n\}$, where $\sigma_i$ represents the set of conditions for the selection of the $S_i$ modes $\sigma_i = \{ condition_1, ..., condition_k\}$. Conditions can be either specific components of the system (e.g., laser working) or measured values of attributes of the system or environment (e.g., battery consumption below $0.5$).

\begin{gather}
    A = \{a_1, ..., a_m\} \\
    F = \{f_1, ..., f_n\} \\
    Modes = \{m_1, ..., m_k\}\\
\end{gather}

\section{Type theoretical formulation}
REPEATED. OLD TEXT NEEDS TO BE RECHECKED
\begin{gather}
    Action: Type \\
    Function: Type \\
    Mode: Type \\
    Component: Type \\
    Condition: Type 
\end{gather}

\begin{equation}
    Monitor(PIType) \rightarrow Measure
\end{equation}

\begin{equation}
    Analyze(m: Measure, ra: RequiredAction) \rightarrow SystemStatus
\end{equation}

\begin{equation}
    Plan(s: SystemStatus) \rightarrow DesiredConfiguration
\end{equation}

\begin{equation}
    Execute(c: DesiredConfiguration) \rightarrow \emptyset
\end{equation}

\begin{equation}
    TaskExeutor(a: Action) \rightarrow RequiredAction
\end{equation}

Architectural adaptation:
\begin{equation}
    architectural(SystemStatus) -> RequiredComponents 
\end{equation}

Parameter adaptation:
\begin{equation}
    parameter(SystemStatus) -> RequiredComponentParameters 
\end{equation}

Task-and-architectural co-adaptation: 
\begin{equation}
    taca(SystemStatus) -> RequiredActions, RequiredConfiguration 
\end{equation}
Also goals as input?

\section{Execute}

\begin{lstlisting}
execute(fds, c_configs):
    c_activate, c_deactivate = get_component_adaptation(fds)
    result_d = deactivate_components(c_deactivate)
    result_a = deactivate_components(c_activate)
    system_modes= get_sytem_modes(fds)
    result_sm = set_system_modes(system_modes)
    result_pa = execute_param_adaptation(c_configs)
    kb = update_kb(kb, rd, ra, sn, rpa)
\end{lstlisting}

\section{ROSA type theoretical formulation}
Based on the PERA type system proposed in~\cite{}

\subsection{ROSA model:}

Entities:
\begin{gather}
    Action: Ent \\
    Function: Ent \\
    Component: Ent \\
    ComponentParameter: Ent \\
    Measure: Ent \\
    QualityAttribute: Ent \\
    QualityAttribute \leq : Measure \\
    EnvironmentalAttribute: Ent \\
    EnvironmentalAttribute \leq : Measure  \\
\end{gather}

Relationships:
\begin{gather}
    functional\_requirement : Obj({action, required\_function}) \\
    action : Itf \\
    required\_function : Itf \\
    function\_design : Obj({function, required\_component}) \\
    function : Itf \\
    required\_component : Itf \\
    component\_configuration : Obj({component, parameter}) \\
    component : Itf \\
    parameter : Itf \\
\end{gather}

Attribute Interfaces:

\begin{gather}
    status : Att(\{status\_owner\}) \\
    status_owner : Itf \\
    status \leq \mathbb{S} \\
    component\_status : Att(\{component\_status\_owner\}) \\
    component\_status_owner : Itf \\
    component\_status \leq status \\
\end{gather}

Interfaces implementation:

\subsection{ROSA loop}

Monitor:
\begin{gather}
    SensorInput : Type \\
    MonitorOutput: Type \\
    measurement \leq MonitorOutput \\
    % cs : component\_status(\{c:component\_status_owner\}) \\
    component\_status \leq MonitorOutput \\
    s: SensorInput \vdash monitor(s) : MonitorOutput \\
\end{gather}

Analyze:
\begin{gather}
    AnalyzeOutput: Type \\
    function\_design\_status \leq AnalyzeOutput \\
    component\_configuration\_status \leq AnalyzeOutput \\
    action\_status \leq AnalyzeOutput \\
    mo : MonitorOutput \vdash analyze(mo) : \{ ao>AnalyzeOutput \}\\
\end{gather}

Should the analyze output be the system status or the functions and components that need to be adapted?


Plan:
\begin{gather}
    % PlanOutput: Type \\
    % reconfiguration\_plan : PlanOutput \\
    ao: AnalyzeOutput \vdash plan(ao) : reconfiguration\_plan 
\end{gather}

Execute:
\begin{gather}
    ExecuteOutput : Type \\
    component\_active \leq ExecuteOutput \\
    component\_parameter\_active \leq ExecuteOutput \\
    component\_process \leq ExecuteOutput \\
    reconfiguration\_plan\_result \leq ExecuteOutput \\
    rp : reconfiguration\_plan \vdash execute(rp) : \{ eo>ExecuteOutput \}, 
\end{gather}

ROSA loop:
\begin{lstlisting}
    while true:
        s = sensor_inputs()
        mo = monitor(s)
        insert_kb(mo)

    while true: # not really a loop
        mo = fetch_kb()
        ao = analyze(mo)
        insert_kb(ao)
    
    while true: # not really a loop
        ao = fetch_kb()
        rp = configuration\_plan(ao)

    while true: # not really a loop
        rp = fetch_kb()
        execute(rp)

    while true: # not really a loop
        ao = fetch_kb()
        task_plan(ao)

\end{lstlisting}